\section{Literature and Contribution}

Research on Chinese local finance explains why LGFVs became central to urban
development. Local governments combined broad expenditure responsibilities,
constrained formal borrowing, land-centered growth, and strong investment
incentives by using local state-owned enterprises and platform companies to
finance infrastructure and development \citep{montinola1995,jin2005,oi1992,
oi1995,wong2013,tsui2011}. Post-2008 studies show how this arrangement shifted
bank lending toward municipal corporate bonds and shadow-banking instruments
and produced durable debt obligations \citep{bai2016longshadow,
chen_he_liu_2020_financing,liu_oi_zhang_2022_grand_bargain}. This literature
establishes the fiscal function of LGFVs, but debt levels and issuer counts do
not identify the institutional location of that function after an exit event.

Recent work comes closer to the paper's measurement problem. Studies of
market-oriented transformation distinguish changes in business, management,
and financing, and they warn that asset integration can be formal without
substantive control \citep{li_feng_hao_2022_lgfp_transformation}. Research on
classified regulation finds that exit from the monitoring list can improve
some performance measures while debt continues to grow
\citep{fang_lu_2024_classified_regulation}. Issuer declarations that a firm no
longer performs government financing can affect borrowing costs even when
financing shifts to other platforms \citep{li_he_ai_2025_transformation}. A
longitudinal study of Jiaxing likewise shows that list exit coexisted with
continuing public ownership, infrastructure functions, government-linked
repayment, and later regrouping \citep{feng_wu_zhang_2022_jiaxing}. These
studies preclude a claim that this paper first noticed the gap between list
status and fiscal function.

The organizational-compliance literature supplies the broader conceptual
language. Policy-practice decoupling separates formal policy from what an
organization does, while means-ends decoupling separates organizational
practice from the outcome that regulation seeks
\citep{bromley_powell_2012_decoupling,de_bree_stoopendaal_2020_recoupling}.
Empirical work also distinguishes legal from practical compliance and shows
that administrative capacity and monitoring can affect that gap
\citep{short_toffel_2010_self_regulation,
zhelyazkova_kaya_schrama_2016_compliance}. LGFV functional transfer is not
identical to either form. The original issuer may genuinely stop performing the
function while another public entity assumes it. The relevant measurement task
is therefore to locate the function, not simply to code whether the original
organization still complies.

The text-measurement literature sets a second boundary. Codebooks, language
models, and high predictive accuracy do not by themselves validate a downstream
substantive estimate. Validation is application-specific, class-specific error
matters, and learned proxies can attenuate, exaggerate, or reverse a regression
coefficient \citep{grimmer_stewart_2013_text,knox_lucas_cho_2022_proxies,
teblunthuis_hase_chan_2024_misclassification,
halterman_keith_2026_codebook}. Design-based supervised learning offers a way
to use imperfect surrogate outcomes when human labels are collected with known
sampling probabilities \citep{egami2023dsl}. The present paper accordingly
treats LLM labels as screening outputs rather than as observed institutional
fates.

The contribution lies in combining four elements at the city-platform level:
a verified formal event, post-event tracing of the financing and project
function, a mutually exclusive four-fate classification, and a retained source
packet that permits review. Within the audited literature, no study applies
that joint design across a multi-case LGFV sample. The claim is deliberately
bounded because Crossref and OpenAlex incompletely index Chinese-language social
science and because an absence in a targeted search cannot establish universal
priority. The framework should be judged by whether it produces transparent,
replicable distinctions that official exit counts cannot make.
